\documentclass[11pt]{article}
\usepackage[a4paper,margin=25mm]{geometry}
\usepackage{amsmath,amssymb,amsthm,booktabs,array}
\usepackage[hidelinks]{hyperref}
\usepackage[T1]{fontenc}
\usepackage{lmodern,microtype}
\newtheorem{theorem}{Theorem}
\newtheorem{lemma}[theorem]{Lemma}
\newtheorem{corollary}[theorem]{Corollary}
\newcommand{\inv}{\operatorname{inv}}
\title{Real positive-definite counterexamples to\newline extended $q$-permanent monotonicity}
\author{Dev D. Goyal}
\date{26 October 2026}
\usepackage[numbers]{natbib}
\usepackage{xurl}
\begin{document}
\maketitle
\begin{abstract}
An explicit real symmetric positive-definite matrix of order five has an inversion-weighted $q$-permanent that decreases between $q=14$ and $q=15$. Exactly nine permutations contribute nonzero terms. This refutes the real positive-definite monotonicity conjecture in da Fonseca's 2010 paper. A sparse family gives a counterexample at every prescribed real $q_0>1$. For each such parameter, a real symmetric positive-definite matrix has a strictly negative $q$-permanent derivative at $q_0$. The dimension may depend on $q_0$. Lean 4 formalizations establish the finite example, the theorem for every prescribed parameter, and the correlation and integer-entry consequences. The result does not address real-symmetric monotonicity on $[-1,1]$.
\end{abstract}

\section{The monotonicity conjecture}
The ordinary permanent sums one entry product for each permutation. The $q$-permanent assigns an additional weight to each product. For an $n\times n$ matrix $A=(a_{ij})$, define
\[
 P_q(A)=\sum_{\sigma\in S_n}q^{\inv(\sigma)}\prod_{i=1}^n a_{i,\sigma(i)},\qquad
 \inv(\sigma)=\#\{(i,j):i<j,\ \sigma(i)>\sigma(j)\}.
\]
The row and column order is fixed. A simultaneous permutation of rows and columns can change this polynomial. The convention is $0^0=1$.

Da Fonseca's 2010 paper uses $A>0$ for a real symmetric positive-definite matrix. Under this convention, Conjecture~1 in Section~2 asserts strict monotonicity on an interval with left endpoint at or below $-1$ and right endpoint $+\infty$ \cite{fonseca2010}. Its later restatement uses a left endpoint strictly below $-1$ \cite[Conjecture~2]{fonseca2018}. Either formulation implies
\begin{equation}\label{eq:implied}
 P_{q_1}(A)<P_{q_2}(A)\qquad(-1<q_1<q_2)
\end{equation}
for every admissible non-diagonal real symmetric positive-definite $A$. The counterexample contradicts \eqref{eq:implied} at two positive integer parameters. It therefore refutes both formulations.

The Bapat--Lal monotonicity assertion is restricted to $[-1,1]$. The October 2026 record of Kitamura's order-144 complex Hermitian counterexample explicitly leaves its real-symmetric case unsettled \cite{kitamura}. The present results concern $q>1$.

Mitchell discusses a failure of the extended conjecture for positive
semidefinite rank-one matrices at negative parameters \cite{mitchell}.
That result does not supply the real positive-definite examples above one
proved here. De S\'a's earlier counterexamples concern incorrect tree
expansion formulas \cite{desa}; the present proof uses a direct permutation
sum and does not use those formulas.

\section{A five-dimensional exact counterexample}
\begin{theorem}\label{thm:five}
The matrix
\[
 A=\begin{pmatrix}
 5&0&0&4&1\\
 0&5&0&-1&4\\
 0&0&1&0&0\\
 4&-1&0&5&0\\
 1&4&0&0&5
 \end{pmatrix}
\]
is real symmetric, non-diagonal, and positive definite. Its $q$-permanent is
\begin{align}\label{eq:five-poly}
 P_q(A)
 &=625+25q^3+800q^5+25q^7+256q^8-32q^9+q^{10}\\
 &=625+25q^3+800q^5+25q^7+q^8(q-16)^2.\nonumber
\end{align}
It satisfies
\[
 P_{14}(A)=8\,968\,822\,249
 >7\,441\,960\,000=P_{15}(A),
\]
with exact difference $1\,526\,862\,249$. Consequently the real monotonicity conjecture described above is false.
\end{theorem}
\begin{proof}
To prove positive definiteness, order the coordinates as $(1,2,4,5,3)$. The resulting matrix is the direct sum of $[1]$ and
\[
 \begin{pmatrix}5I_2&R\\R^T&5I_2\end{pmatrix},\qquad
 R=\begin{pmatrix}4&1\\-1&4\end{pmatrix},\qquad RR^T=17I_2.
\]
Its eigenvalues are $1$ and $5\pm\sqrt{17}$, with each of the latter occurring twice. All are positive. For the $q$-permanent calculation, retain the displayed order of the matrix in the theorem.

Every nonzero term fixes coordinate $3$. Among the remaining coordinates, four transpositions correspond to nonzero off-diagonal entries. There are also two disjoint pairs of these transpositions and two oriented four-cycles. Together with the identity, these give exactly the nine terms in Table~\ref{tab:terms}. Summing the table proves \eqref{eq:five-poly}. Substitution of $14$ and $15$ yields the stated integer values, which contradict strict monotonicity.
\end{proof}

\begin{table}[htbp]
\centering
\begin{tabular}{@{}lrr@{}}
\toprule
Permutation in one-line notation & Inversions & Entry product\\
\midrule
$(1,2,3,4,5)$ &0&625\\
$(1,4,3,2,5)$ &3&25\\
$(1,5,3,4,2)$ &5&400\\
$(4,2,3,1,5)$ &5&400\\
$(4,5,3,1,2)$ &8&256\\
$(4,5,3,2,1)$ &9&$-16$\\
$(5,2,3,4,1)$ &7&25\\
$(5,4,3,1,2)$ &9&$-16$\\
$(5,4,3,2,1)$ &10&1\\
\bottomrule
\end{tabular}
\caption{All nonzero terms for the fixed order in Theorem~\ref{thm:five}.}
\label{tab:terms}
\end{table}

\section{A sparse family}
Let $m\ge0$ be an integer and $n=m+4$. The active coordinates are
\[
 i=1,\qquad j=2,\qquad k=m+3,\qquad \ell=m+4.
\]
For real $a\ge0$ and $d>\sqrt{a^2+1}$, define $A_m(a,d)$ as follows. It equals the identity on the $m$ remaining coordinates. All entries between these coordinates and the active coordinates are zero. On the active coordinates, in increasing order, its principal submatrix is
\begin{equation}\label{eq:core}
 \begin{pmatrix}
 d&0&a&1\\
 0&d&-1&a\\
 a&-1&d&0\\
 1&a&0&d
 \end{pmatrix}.
\end{equation}
It is positive definite, since the active block has eigenvalues $d\pm\sqrt{a^2+1}$, each twice. The intervening identity coordinates are part of the fixed ordering and must not be deleted when counting inversions.

\begin{lemma}\label{lem:family}
For every $m\ge0$,
\begin{equation}\label{eq:family}
 P_q(A_m(a,d))
 =d^4+d^2q^{2m+1}(1+2a^2q^2+q^4)
   +q^{4m+4}(a^2-q)^2.
\end{equation}
\end{lemma}
\begin{proof}
The isolated coordinates are fixed in every nonzero term. The following table uses cycle notation for the nine possible active permutations. Inversion counts use the full order $1,\ldots,m+4$.
\[
\begin{array}{c|c|c}
\text{Active permutation}&\text{Inversions}&\text{Entry product}\\\hline
\mathrm{id}&0&d^4\\
(j\ k)&2m+1&d^2\\
(i\ k),\ (j\ \ell)&2m+3&a^2d^2\ \text{each}\\
(i\ \ell)&2m+5&d^2\\
(i\ k)(j\ \ell)&4m+4&a^4\\
(i\ k\ j\ \ell),\ (i\ \ell\ j\ k)&4m+5&-a^2\ \text{each}\\
(i\ \ell)(j\ k)&4m+6&1
\end{array}
\]
For a transposition of positions $r<s$, the inversion count is $2(s-r)-1$. For each of the last three types, the four active coordinates contribute respectively $4$, $5$, and $6$ inversions. Each intervening fixed coordinate contributes four additional inversions. These facts prove the stated counts. Summing gives
\[
 d^4+d^2q^{2m+1}+2a^2d^2q^{2m+3}+d^2q^{2m+5}
 +a^4q^{4m+4}-2a^2q^{4m+5}+q^{4m+6},
\]
which is \eqref{eq:family}.
\end{proof}

\section{Failure at every prescribed parameter above one}
The last term in \eqref{eq:family} can have a negative derivative.
The proof below chooses $a$ and $d$ in terms of $m$ and a fixed $r>1$.
At $q=r$, this negative contribution grows faster in magnitude with $m$
than the positive derivative of the middle term.
\begin{theorem}\label{thm:allq}
For every real $r>1$, there are an integer $n\ge5$ and a real symmetric positive-definite matrix $A$ of order $n$ such that
\[
 \left.\frac{d}{dq}P_q(A)\right|_{q=r}<0.
\]
Thus no interval $[1,1+\delta]$ with $\delta>0$ can enjoy $q$-permanent monotonicity uniformly over all real symmetric positive-definite matrices of all orders.
\end{theorem}
\begin{proof}
Fix $r>1$. Choose an integer $m\ge1$ sufficiently large that
\begin{equation}\label{eq:m-bound}
 r^{2m+5}>
 (4m+4)(2m+5)\left(\frac{5r}{4}+2\right)
 \left(1+\frac52r^3+r^4\right).
\end{equation}
Such an integer exists. The left side grows exponentially in $m$. The right side is a quadratic polynomial in $m$ times a constant that depends only on $r$.

Set
\[
 K=4m+4,\qquad H=2m+1,\qquad
 t=r\left(1+\frac1K\right),\qquad
 a=\sqrt{t},\qquad d=\sqrt{t+2}.
\]
Keep these parameters fixed when differentiating with respect to $q$. Then $d>\sqrt{a^2+1}$, so $A=A_m(a,d)$ is positive definite. By Lemma~\ref{lem:family},
\[
 P_q(A)=d^4+d^2q^H(1+2tq^2+q^4)+q^K(t-q)^2.
\]
The derivative of the last summand at $q=r$ is exactly
\begin{align*}
 Kr^{K-1}(t-r)^2-2r^K(t-r)
 &=Kr^{K-1}\frac{r^2}{K^2}-2r^K\frac rK\\
 &=-\frac{r^{K+1}}K.
\end{align*}
The derivative of the middle summand is
\[
 (t+2)r^{H-1}\bigl(H+2t(H+2)r^2+(H+4)r^4\bigr).
\]
Since $K\ge4$, the bound $t\le5r/4$ holds. Also $H,H+2\le H+4=2m+5$. Consequently this positive derivative is at most
\[
 \left(\frac{5r}{4}+2\right)r^{H-1}(2m+5)
 \left(1+\frac52r^3+r^4\right).
\]
Multiply \eqref{eq:m-bound} by $r^{H-1}/K$. The result shows that this upper bound is strictly smaller than $r^{K+1}/K$, since $K-H+2=2m+5$. The total derivative is therefore negative.

Finally, given $\delta>0$, apply the theorem with $r=1+\delta/2$. A negative derivative of a polynomial gives a strictly decreasing pair of sufficiently nearby parameters contained in $[1,1+\delta]$.
\end{proof}

\begin{corollary}\label{cor:correlations}
The matrices in Theorems~\ref{thm:five} and \ref{thm:allq} can be replaced by real positive-definite correlation matrices, with the same strict monotonicity failures. At every prescribed $r>1$, a counterexample can additionally be chosen with no zero entries.
\end{corollary}
\begin{proof}
If $D=\operatorname{diag}(s_1,\ldots,s_n)$ with all $s_i>0$, then term by term
\[
 P_q(DAD)=\left(\prod_{i=1}^n s_i^2\right)P_q(A).
\]
Choose $s_i=a_{ii}^{-1/2}$ to normalize the diagonal. Positive definiteness is preserved, and the multiplicative factor is positive and independent of $q$.

For the final assertion, first replace a matrix from Theorem~\ref{thm:allq} by $A+\varepsilon\mathbf1\mathbf1^T$. This remains positive definite for every $\varepsilon>0$. For sufficiently small positive $\varepsilon$, all entries are nonzero. The derivative remains negative because it is a polynomial in the entries and is therefore continuous. Then normalize the diagonal as above.
\end{proof}

\begin{corollary}\label{cor:integer}
For every real $r>1$, a real symmetric positive-definite matrix with integer entries can be chosen such that $P'_r(A)<0$.
\end{corollary}
\begin{proof}
Positive definiteness and the strict inequality $P'_r(A)<0$ are open conditions in the real vector space of symmetric matrices. Approximate the matrix of Theorem~\ref{thm:allq} by a sufficiently close rational symmetric matrix $B$ preserving both conditions. Multiplication by a positive common denominator $c$ gives an integer matrix, and $P'_r(cB)=c^nP'_r(B)<0$.
\end{proof}

\section{Formal verification and independent checks}
The Lean formalization is in
\path{lean/MathAndTheory/Mathematics/QPermanent/}.
Its namespace is \texttt{MathAndTheory.QPermanent}.
The table maps the paper's results to Lean declarations.
\begin{center}\small
\begin{tabular}{@{}p{0.30\linewidth}p{0.65\linewidth}@{}}\toprule
Paper statement & Lean declaration\\\midrule
Theorem~\ref{thm:five}, positivity & \path{counterexample_posDef}\\
Theorem~\ref{thm:five}, polynomial & \path{counterexample_polynomial}\\
Source conjecture is false & \path{not_daFonseca_extended_conjecture}\\
Lemma~\ref{lem:family} & \path{qPermanent_family_formula}\\
Theorem~\ref{thm:allq} & \path{exists_posDef_negative_qPermanent_derivative_ge_five}\\
Correlation consequence & \path{exists_correlation_counterexample}\\
Dense correlation consequence & \path{exists_dense_correlation_negative_qPermanent_derivative}\\
Integer-entry consequence & \path{exists_integer_posDef_negative_qPermanent_derivative}\\\bottomrule
\end{tabular}
\end{center}
The formal definition sums over every permutation and counts ordinary
inversions in the displayed coordinate order. The proof establishes the
contribution of every intervening coordinate. Positive definiteness is
proved by quadratic forms. The integer-entry proof uses rational
approximations to the explicit family parameters, followed by denominator
clearing. This gives the same existence conclusion as
Corollary~\ref{cor:integer}.

The independent program reconstructs the five-dimensional polynomial by
enumerating all $120$ permutations. It checks the positive rational $LDL^T$ pivots
$5,5,1,8/5,8/5$, both endpoint values, and
$P'_{15}(A)=-1\,563\,030\,000$.
For $m=0,1,2,3,4$, it independently enumerates all permutations of the
corresponding family matrix. It also checks the growth condition and
negative derivative with exact rational arithmetic at
$r=101/100,11/10,3/2,2,15$.
The respective choices of $m$ are $925,69,12,6,1$.
These computations check finitely many examples. Theorem~\ref{thm:allq}
and its formalization prove the result for every $r$.

The repository pins Lean 4.35.0 and Mathlib commit
\path{56a005855ea17e0b0958f73ec801ae140b0a1679}.
See \path{lean/docs/q-permanent.md} for statement alignment and
\path{verification/q-permanent/} for commands, software versions, and
recorded results. The Challenge file states the permutation definition
directly and does not import the proof modules. Comparator checks that the
Solution file proves the same statements.

\section{Scope and prior work}
The construction bounds the dimension of a counterexample. It does not
determine the smallest possible dimension. The integer-entry and correlation
conclusions are separate. They do not assert an integer correlation matrix. The results do not settle the
real Bapat--Lal assertion on $[-1,1]$.

No claim of first priority is made.

\section*{Autonomous AI research}
This work is fully autonomous agentic research. The initial research and proofs were generated autonomously by GPT 6 with Pro Reasoning in ChatGPT. GPT 6 Astra in Codex autonomously produced the formalizations and corrections. The authors checked the proofs to the best of their mathematical abilities. This check is not expert validation or peer review. Dev D. Goyal is responsible for publication. This preprint has not been peer reviewed.

\bibliographystyle{plainnat}
\bibliography{references}
\end{document}
